
\documentclass{article}
\usepackage{colortbl}
\usepackage{makecell}
\usepackage{multirow}
\usepackage{supertabular}

\begin{document}

\newcounter{utterance}

\centering \large Interaction Transcript for game `function\_detective', experiment `two\_numbers\_passive\_examples', episode 259 with Qwen3.6{-}35B{-}A3B{-}FP8{-}without{-}reasoning.
\vspace{24pt}

{ \footnotesize  \setcounter{utterance}{1}
\setlength{\tabcolsep}{0pt}
\begin{supertabular}{c@{$\;$}|p{.15\linewidth}@{}p{.15\linewidth}p{.15\linewidth}p{.15\linewidth}p{.15\linewidth}p{.15\linewidth}}
    \# & \multicolumn{2}{c}{Player} && \multicolumn{2}{c}{Game Master} \\
    \hline

    \theutterance \stepcounter{utterance}  
    & & & \multicolumn{4}{p{0.6\linewidth}}{
        \cellcolor[rgb]{0.9,0.9,0.9}{
            \makecell[{{p{\linewidth}}}]{
                \texttt{\tiny{[P1$\langle$GM]}}
                \texttt{I have a hidden Python function for you to figure out.} \\
\\ 
\texttt{Your job is to infer what it does and then submit a correct Python implementation.} \\
\\ 
\texttt{Function signature:} \\
\texttt{(a: int, b: int) {-}> int} \\
\\ 
\texttt{You have 15 turns total.} \\
\\ 
\texttt{I control which valid input{-}output examples you see. You can request them one at a time.} \\
\\ 
\texttt{You must make your best SOLVE attempt no later than the second{-}to{-}last turn.} \\
\texttt{Do not wait until the very end.} \\
\\ 
\texttt{You can use these actions:} \\
\texttt{{-} JSON format described at the end of the prompt} \\
\texttt{{-} SOLVE: ```python ... ```} \\
\\ 
\texttt{Do not add import statements in your guess.} \\
\texttt{Anything you might need is already available, so just write the function body in hypothesis and return the answer.} \\
\\ 
\texttt{Preloaded examples:} \\
\texttt{1. ((0, {-}2), {-}2)} \\
\\ 
\texttt{Scoring:} \\
\texttt{{-} correctness is required for any non{-}zero score} \\
\texttt{{-} in interactive modes, solving with fewer information requests is better} \\
\\ 
\texttt{Input formatting reminders:} \\
\texttt{{-} strings must be quoted} \\
\texttt{{-} booleans must be True/False} \\
\texttt{{-} use commas between multiple values} \\
\texttt{{-} do not write any extra action tags in the same message} \\
\\ 
\texttt{Format example:} \\
\texttt{You write:} \\
\texttt{\{'current\_mode': <probing/confirming/don't know; choose whichever best describes your current state in the task>, 'input\_rationale': <brief description of why you chose this hypothesis>, 'hypothesis': <either the exact string "unknown" if you have no guess yet OR your current best guess as runnable Python matching the given signature, with a one{-}line docstring stating what it does. Never write placeholder code, if you cannot commit to a guess, use "unknown". Match the given signature.} \\
\texttt{e.g.} \\
\texttt{def f(x):} \\
\texttt{	"""Returns the absolute value of x."""} \\
\texttt{	return x if x >= 0 else {-}x} \\
\\ 
\texttt{Use "unknown" if you have no guess.>, 'confidence\_score': <score between 1{-}100. 100 means highest confidence in your hypothesis>\}} \\
\\ 
\texttt{Output only one valid json per turn.} \\
            }
        }
    }
     \\ \\

    \theutterance \stepcounter{utterance}  
    & \multicolumn{4}{p{0.6\linewidth}}{
        \cellcolor[rgb]{0.561,0.737,0.561}{
            \makecell[{{p{\linewidth}}}]{
                \texttt{\tiny{[P1$\rangle$GM]}}
                \texttt{\{'current\_mode': 'probing', 'input\_rationale': 'I have only one example so far: f(0, {-}2) = {-}2. This is consistent with several simple functions like returning the second argument, returning the minimum, or returning the sum. I need more data points to distinguish between these possibilities. I will request an example with positive integers to see if the function behaves differently.', 'hypothesis': 'unknown', 'confidence\_score': 10\}} \\
            }
        }
    }
    & & \\ \\

    \theutterance \stepcounter{utterance}  
    & & & \multicolumn{4}{p{0.6\linewidth}}{
        \cellcolor[rgb]{0.9,0.9,0.9}{
            \makecell[{{p{\linewidth}}}]{
                \texttt{\tiny{[P1$\langle$GM]}}
                \texttt{Game Violation: ParseError: Response is not valid JSON.} \\
\texttt{Output ONLY:} \\
\texttt{One JSON with confidence\_score, current\_mode, input\_rationale, and hypothesis keys ...} \\
\texttt{OR} \\
\texttt{SOLVE: ```python ... ```} \\
            }
        }
    }
     \\ \\

    \theutterance \stepcounter{utterance}  
    & \multicolumn{4}{p{0.6\linewidth}}{
        \cellcolor[rgb]{0.561,0.737,0.561}{
            \makecell[{{p{\linewidth}}}]{
                \texttt{\tiny{[P1$\rangle$GM]}}
                \texttt{\{'current\_mode': 'probing', 'input\_rationale': 'I have only one example so far: f(0, {-}2) = {-}2. This is consistent with several simple functions like returning the second argument, returning the minimum, or returning the sum. I need more data points to distinguish between these possibilities. I will request an example with positive integers to see if the function behaves differently.', 'hypothesis': 'unknown', 'confidence\_score': 10\}} \\
            }
        }
    }
    & & \\ \\

    \theutterance \stepcounter{utterance}  
    & & & \multicolumn{4}{p{0.6\linewidth}}{
        \cellcolor[rgb]{0.9,0.9,0.9}{
            \makecell[{{p{\linewidth}}}]{
                \texttt{\tiny{[P1$\langle$GM]}}
                \texttt{Game Violation: ParseError: Response is not valid JSON.} \\
\texttt{Output ONLY:} \\
\texttt{One JSON with confidence\_score, current\_mode, input\_rationale, and hypothesis keys ...} \\
\texttt{OR} \\
\texttt{SOLVE: ```python ... ```} \\
            }
        }
    }
     \\ \\

    \theutterance \stepcounter{utterance}  
    & \multicolumn{4}{p{0.6\linewidth}}{
        \cellcolor[rgb]{0.561,0.737,0.561}{
            \makecell[{{p{\linewidth}}}]{
                \texttt{\tiny{[P1$\rangle$GM]}}
                \texttt{\{"current\_mode": "probing", "input\_rationale": "I have only one example: f(0, {-}2) = {-}2. This could be min(a,b), max(a,b), a+b, b, etc. I need more examples to distinguish. I will request f(1, 2) to see if it returns 1, 2, or 3.", "hypothesis": "unknown", "confidence\_score": 10\}} \\
            }
        }
    }
    & & \\ \\

    \theutterance \stepcounter{utterance}  
    & & & \multicolumn{4}{p{0.6\linewidth}}{
        \cellcolor[rgb]{0.9,0.9,0.9}{
            \makecell[{{p{\linewidth}}}]{
                \texttt{\tiny{[P1$\langle$GM]}}
                \texttt{OUTPUT: (({-}1, 99), {-}1)} \\
            }
        }
    }
     \\ \\

    \theutterance \stepcounter{utterance}  
    & \multicolumn{4}{p{0.6\linewidth}}{
        \cellcolor[rgb]{0.561,0.737,0.561}{
            \makecell[{{p{\linewidth}}}]{
                \texttt{\tiny{[P1$\rangle$GM]}}
                \texttt{\{"current\_mode": "probing", "input\_rationale": "I have two examples: f(0, {-}2) = {-}2 and f({-}1, 99) = {-}1. Both results are the minimum of the two inputs. This strongly suggests the function returns min(a, b). I will request one more example with both positive numbers to confirm.", "hypothesis": "unknown", "confidence\_score": 60\}} \\
            }
        }
    }
    & & \\ \\

    \theutterance \stepcounter{utterance}  
    & & & \multicolumn{4}{p{0.6\linewidth}}{
        \cellcolor[rgb]{0.9,0.9,0.9}{
            \makecell[{{p{\linewidth}}}]{
                \texttt{\tiny{[P1$\langle$GM]}}
                \texttt{OUTPUT: ((10, {-}1000), {-}1000)} \\
            }
        }
    }
     \\ \\

    \theutterance \stepcounter{utterance}  
    & \multicolumn{4}{p{0.6\linewidth}}{
        \cellcolor[rgb]{0.561,0.737,0.561}{
            \makecell[{{p{\linewidth}}}]{
                \texttt{\tiny{[P1$\rangle$GM]}}
                \texttt{\{"current\_mode": "confirming", "input\_rationale": "I have three examples: f(0, {-}2) = {-}2, f({-}1, 99) = {-}1, and f(10, {-}1000) = {-}1000. In all cases, the result is the minimum of the two inputs. This is a very strong pattern. I am confident the function is min(a, b).", "hypothesis": "def f(a: int, b: int) {-}> int:\textbackslash{}n    \textbackslash{}"\textbackslash{}"\textbackslash{}"Returns the minimum of two integers.\textbackslash{}"\textbackslash{}"\textbackslash{}"\textbackslash{}n    return a if a < b else b", "confidence\_score": 95\}} \\
            }
        }
    }
    & & \\ \\

    \theutterance \stepcounter{utterance}  
    & & & \multicolumn{4}{p{0.6\linewidth}}{
        \cellcolor[rgb]{0.9,0.9,0.9}{
            \makecell[{{p{\linewidth}}}]{
                \texttt{\tiny{[P1$\langle$GM]}}
                \texttt{OUTPUT: ((1, 99), 0)} \\
            }
        }
    }
     \\ \\

    \theutterance \stepcounter{utterance}  
    & \multicolumn{4}{p{0.6\linewidth}}{
        \cellcolor[rgb]{0.561,0.737,0.561}{
            \makecell[{{p{\linewidth}}}]{
                \texttt{\tiny{[P1$\rangle$GM]}}
                \texttt{\{"current\_mode": "probing", "input\_rationale": "The previous hypothesis (min(a, b)) is contradicted by the new example: f(1, 99) = 0, whereas min(1, 99) = 1. The result 0 is neither a, b, min, max, nor sum. It looks like it might be related to bitwise operations or some other logic. Let's look at the examples again: (0, {-}2) {-}> {-}2, ({-}1, 99) {-}> {-}1, (10, {-}1000) {-}> {-}1000, (1, 99) {-}> 0. The first three were min. The last one is not. This suggests the function might be min(a, b) but with some condition, or perhaps it's something like a \& b (bitwise AND)? Let's check: 0 \& {-}2 = 0 (no, result is {-}2). {-}1 \& 99 = 1 (no, result is {-}1). 10 \& {-}1000 = 0 (no, result is {-}1000). 1 \& 99 = 1 (no, result is 0). So it's not bitwise AND. What about a \^{} b (XOR)? 0 \^{} {-}2 = {-}2. {-}1 \^{} 99 = 98 (} \\
            }
        }
    }
    & & \\ \\

    \theutterance \stepcounter{utterance}  
    & & & \multicolumn{4}{p{0.6\linewidth}}{
        \cellcolor[rgb]{0.9,0.9,0.9}{
            \makecell[{{p{\linewidth}}}]{
                \texttt{\tiny{[P1$\langle$GM]}}
                \texttt{Game Violation: ParseError: Response is not valid JSON.} \\
\texttt{Output ONLY:} \\
\texttt{One JSON with confidence\_score, current\_mode, input\_rationale, and hypothesis keys ...} \\
\texttt{OR} \\
\texttt{SOLVE: ```python ... ```} \\
            }
        }
    }
     \\ \\

    \theutterance \stepcounter{utterance}  
    & \multicolumn{4}{p{0.6\linewidth}}{
        \cellcolor[rgb]{0.561,0.737,0.561}{
            \makecell[{{p{\linewidth}}}]{
                \texttt{\tiny{[P1$\rangle$GM]}}
                \texttt{\{"current\_mode": "probing", "input\_rationale": "The hypothesis min(a,b) is disproven by f(1,99)=0. I need more examples to find the pattern. I will test f(2, 3) to see if it returns 0 or something else.", "hypothesis": "unknown", "confidence\_score": 20\}} \\
            }
        }
    }
    & & \\ \\

    \theutterance \stepcounter{utterance}  
    & & & \multicolumn{4}{p{0.6\linewidth}}{
        \cellcolor[rgb]{0.9,0.9,0.9}{
            \makecell[{{p{\linewidth}}}]{
                \texttt{\tiny{[P1$\langle$GM]}}
                \texttt{OUTPUT: ((1000, 99), 0)} \\
            }
        }
    }
     \\ \\

    \theutterance \stepcounter{utterance}  
    & \multicolumn{4}{p{0.6\linewidth}}{
        \cellcolor[rgb]{0.561,0.737,0.561}{
            \makecell[{{p{\linewidth}}}]{
                \texttt{\tiny{[P1$\rangle$GM]}}
                \texttt{\{"current\_mode": "probing", "input\_rationale": "I have four examples: f(0, {-}2) = {-}2, f({-}1, 99) = {-}1, f(10, {-}1000) = {-}1000, f(1, 99) = 0, f(1000, 99) = 0. The first three look like min(a,b). The last two are 0. Notice that in the last two cases, both inputs are positive. In the first three, at least one input is non{-}positive (0 or negative). This suggests the function might be: if a > 0 and b > 0, return 0; else return min(a, b). Let's test this hypothesis with another case where one is positive and one is negative, or both negative.", "hypothesis": "unknown", "confidence\_score": 40\}} \\
            }
        }
    }
    & & \\ \\

    \theutterance \stepcounter{utterance}  
    & & & \multicolumn{4}{p{0.6\linewidth}}{
        \cellcolor[rgb]{0.9,0.9,0.9}{
            \makecell[{{p{\linewidth}}}]{
                \texttt{\tiny{[P1$\langle$GM]}}
                \texttt{OUTPUT: ((1000, 0), 0)} \\
            }
        }
    }
     \\ \\

    \theutterance \stepcounter{utterance}  
    & \multicolumn{4}{p{0.6\linewidth}}{
        \cellcolor[rgb]{0.561,0.737,0.561}{
            \makecell[{{p{\linewidth}}}]{
                \texttt{\tiny{[P1$\rangle$GM]}}
                \texttt{\{"current\_mode": "probing", "input\_rationale": "I have five examples: f(0, {-}2) = {-}2, f({-}1, 99) = {-}1, f(10, {-}1000) = {-}1000, f(1, 99) = 0, f(1000, 99) = 0, f(1000, 0) = 0. The pattern seems to be: if both a and b are non{-}negative (>= 0), return 0. Otherwise, return min(a, b). Let's verify: (0, {-}2): 0>=0 but {-}2<0, so not both non{-}negative. min(0, {-}2) = {-}2. Correct. ({-}1, 99): {-}1<0, so not both non{-}negative. min({-}1, 99) = {-}1. Correct. (10, {-}1000): {-}1000<0, so not both non{-}negative. min(10, {-}1000) = {-}1000. Correct. (1, 99): both >=0. Return 0. Correct. (1000, 99): both >=0. Return 0. Correct. (1000, 0): both >=0. Return 0} \\
            }
        }
    }
    & & \\ \\

    \theutterance \stepcounter{utterance}  
    & & & \multicolumn{4}{p{0.6\linewidth}}{
        \cellcolor[rgb]{0.9,0.9,0.9}{
            \makecell[{{p{\linewidth}}}]{
                \texttt{\tiny{[P1$\langle$GM]}}
                \texttt{Game Violation: ParseError: Response is not valid JSON.} \\
\texttt{Output ONLY:} \\
\texttt{One JSON with confidence\_score, current\_mode, input\_rationale, and hypothesis keys ...} \\
\texttt{OR} \\
\texttt{SOLVE: ```python ... ```} \\
            }
        }
    }
     \\ \\

    \theutterance \stepcounter{utterance}  
    & \multicolumn{4}{p{0.6\linewidth}}{
        \cellcolor[rgb]{0.561,0.737,0.561}{
            \makecell[{{p{\linewidth}}}]{
                \texttt{\tiny{[P1$\rangle$GM]}}
                \texttt{\{"current\_mode": "confirming", "input\_rationale": "I have six examples: f(0, {-}2) = {-}2, f({-}1, 99) = {-}1, f(10, {-}1000) = {-}1000, f(1, 99) = 0, f(1000, 99) = 0, f(1000, 0) = 0. The pattern is consistent: if both a and b are non{-}negative (>= 0), return 0. Otherwise, return min(a, b). I am confident in this hypothesis.", "hypothesis": "def f(a: int, b: int) {-}> int:\textbackslash{}n    \textbackslash{}"\textbackslash{}"\textbackslash{}"Returns 0 if both inputs are non{-}negative, otherwise returns the minimum of the two inputs.\textbackslash{}"\textbackslash{}"\textbackslash{}"\textbackslash{}n    if a >= 0 and b >= 0:\textbackslash{}n        return 0\textbackslash{}n    return a if a < b else b", "confidence\_score": 95\}} \\
            }
        }
    }
    & & \\ \\

    \theutterance \stepcounter{utterance}  
    & & & \multicolumn{4}{p{0.6\linewidth}}{
        \cellcolor[rgb]{0.9,0.9,0.9}{
            \makecell[{{p{\linewidth}}}]{
                \texttt{\tiny{[P1$\langle$GM]}}
                \texttt{OUTPUT: (({-}99, {-}10), {-}109)} \\
            }
        }
    }
     \\ \\

    \theutterance \stepcounter{utterance}  
    & \multicolumn{4}{p{0.6\linewidth}}{
        \cellcolor[rgb]{0.561,0.737,0.561}{
            \makecell[{{p{\linewidth}}}]{
                \texttt{\tiny{[P1$\rangle$GM]}}
                \texttt{\{"current\_mode": "probing", "input\_rationale": "The new example f({-}99, {-}10) = {-}109 contradicts my previous hypothesis. My hypothesis predicted min({-}99, {-}10) = {-}99. The actual result is {-}109, which is the sum of {-}99 and {-}10. Let's re{-}evaluate all examples with the sum hypothesis: f(0, {-}2) = {-}2 (sum). f({-}1, 99) = 98 (sum), but output is {-}1. So it's not just sum. Let's look at the outputs again: {-}2, {-}1, {-}1000, 0, 0, 0, {-}109. Inputs: (0, {-}2), ({-}1, 99), (10, {-}1000), (1, 99), (1000, 99), (1000, 0), ({-}99, {-}10). Notice that when both are positive, output is 0. When one is negative and one positive, output is the negative one (min). When both are negative, output is sum. Let's check: (0, {-}2): 0 is non{-}negative, {-}2 is negative. Not both positive. Not both negative. One neg, one non{-}neg. Output {-}2 (} \\
            }
        }
    }
    & & \\ \\

    \theutterance \stepcounter{utterance}  
    & & & \multicolumn{4}{p{0.6\linewidth}}{
        \cellcolor[rgb]{0.9,0.9,0.9}{
            \makecell[{{p{\linewidth}}}]{
                \texttt{\tiny{[P1$\langle$GM]}}
                \texttt{Game Violation: ParseError: Response is not valid JSON.} \\
\texttt{Output ONLY:} \\
\texttt{One JSON with confidence\_score, current\_mode, input\_rationale, and hypothesis keys ...} \\
\texttt{OR} \\
\texttt{SOLVE: ```python ... ```} \\
            }
        }
    }
     \\ \\

    \theutterance \stepcounter{utterance}  
    & \multicolumn{4}{p{0.6\linewidth}}{
        \cellcolor[rgb]{0.561,0.737,0.561}{
            \makecell[{{p{\linewidth}}}]{
                \texttt{\tiny{[P1$\rangle$GM]}}
                \texttt{\{"current\_mode": "probing", "input\_rationale": "The hypothesis that f(a,b) is 0 if both >=0, else min(a,b) is disproven by f({-}99, {-}10) = {-}109. The result {-}109 is the sum of {-}99 and {-}10. Let's check if the function is: if a >= 0 and b >= 0: return 0; elif a < 0 and b < 0: return a + b; else: return min(a, b). Check examples: (0, {-}2): 0>=0, {-}2<0. Else branch. min(0, {-}2) = {-}2. Correct. ({-}1, 99): {-}1<0, 99>=0. Else branch. min({-}1, 99) = {-}1. Correct. (10, {-}1000): 10>=0, {-}1000<0. Else branch. min(10, {-}1000) = {-}1000. Correct. (1, 99): both >=0. Return 0. Correct. (1000, 99): both >=0. Return 0. Correct. (1000, 0): both >=0. Return 0. Correct. ({-}99, {-}10): both} \\
            }
        }
    }
    & & \\ \\

    \theutterance \stepcounter{utterance}  
    & & & \multicolumn{4}{p{0.6\linewidth}}{
        \cellcolor[rgb]{0.9,0.9,0.9}{
            \makecell[{{p{\linewidth}}}]{
                \texttt{\tiny{[P1$\langle$GM]}}
                \texttt{Game Violation: ParseError: Response is not valid JSON.} \\
\texttt{Output ONLY:} \\
\texttt{One JSON with confidence\_score, current\_mode, input\_rationale, and hypothesis keys ...} \\
\texttt{OR} \\
\texttt{SOLVE: ```python ... ```} \\
            }
        }
    }
     \\ \\

    \theutterance \stepcounter{utterance}  
    & \multicolumn{4}{p{0.6\linewidth}}{
        \cellcolor[rgb]{0.561,0.737,0.561}{
            \makecell[{{p{\linewidth}}}]{
                \texttt{\tiny{[P1$\rangle$GM]}}
                \texttt{\{"current\_mode": "confirming", "input\_rationale": "I have seven examples. The pattern appears to be: if both a and b are non{-}negative, return 0. If both are negative, return their sum. If one is non{-}negative and the other is negative, return the negative one (which is min(a,b)). This covers all cases: (0, {-}2) {-}> {-}2 (mixed, min), ({-}1, 99) {-}> {-}1 (mixed, min), (10, {-}1000) {-}> {-}1000 (mixed, min), (1, 99) {-}> 0 (both non{-}neg), (1000, 99) {-}> 0 (both non{-}neg), (1000, 0) {-}> 0 (both non{-}neg), ({-}99, {-}10) {-}> {-}109 (both neg, sum). This hypothesis fits all data points perfectly.", "hypothesis": "def f(a: int, b: int) {-}> int:\textbackslash{}n    \textbackslash{}"\textbackslash{}"\textbackslash{}"Returns 0 if both inputs are non{-}negative, the sum if both are negative, and the minimum (which will be the negative one) if signs differ.\textbackslash{}"\textbackslash{}"\textbackslash{}"\textbackslash{}n    if a >= 0 and b >= 0:\textbackslash{}n        return 0\textbackslash{}n    elif a < 0 and b < 0:\textbackslash{}n        return a + b\textbackslash{}n} \\
            }
        }
    }
    & & \\ \\

    \theutterance \stepcounter{utterance}  
    & & & \multicolumn{4}{p{0.6\linewidth}}{
        \cellcolor[rgb]{0.9,0.9,0.9}{
            \makecell[{{p{\linewidth}}}]{
                \texttt{\tiny{[P1$\langle$GM]}}
                \texttt{Game Violation: ParseError: Response is not valid JSON.} \\
\texttt{Output ONLY:} \\
\texttt{One JSON with confidence\_score, current\_mode, input\_rationale, and hypothesis keys ...} \\
\texttt{OR} \\
\texttt{SOLVE: ```python ... ```} \\
            }
        }
    }
     \\ \\

    \theutterance \stepcounter{utterance}  
    & \multicolumn{4}{p{0.6\linewidth}}{
        \cellcolor[rgb]{0.561,0.737,0.561}{
            \makecell[{{p{\linewidth}}}]{
                \texttt{\tiny{[P1$\rangle$GM]}}
                \texttt{\{"current\_mode": "confirming", "input\_rationale": "I have seven examples that fit the hypothesis: if both a and b are non{-}negative, return 0; if both are negative, return a + b; otherwise return min(a, b). This covers all observed cases correctly. I am confident in this solution.", "hypothesis": "def f(a: int, b: int) {-}> int:\textbackslash{}n    \textbackslash{}"\textbackslash{}"\textbackslash{}"Returns 0 if both inputs are non{-}negative, the sum if both are negative, and the minimum if signs differ.\textbackslash{}"\textbackslash{}"\textbackslash{}"\textbackslash{}n    if a >= 0 and b >= 0:\textbackslash{}n        return 0\textbackslash{}n    elif a < 0 and b < 0:\textbackslash{}n        return a + b\textbackslash{}n    else:\textbackslash{}n        return a if a < b else b", "confidence\_score": 95\}} \\
            }
        }
    }
    & & \\ \\

    \theutterance \stepcounter{utterance}  
    & & & \multicolumn{4}{p{0.6\linewidth}}{
        \cellcolor[rgb]{0.9,0.9,0.9}{
            \makecell[{{p{\linewidth}}}]{
                \texttt{\tiny{[P1$\langle$GM]}}
                \texttt{OUTPUT: ((1000, 10), 0)} \\
            }
        }
    }
     \\ \\

    \theutterance \stepcounter{utterance}  
    & \multicolumn{4}{p{0.6\linewidth}}{
        \cellcolor[rgb]{0.561,0.737,0.561}{
            \makecell[{{p{\linewidth}}}]{
                \texttt{\tiny{[P1$\rangle$GM]}}
                \texttt{\{"current\_mode": "confirming", "input\_rationale": "The new example f(1000, 10) = 0 is consistent with my hypothesis: both inputs are non{-}negative, so the function returns 0. All 8 examples now fit the hypothesis perfectly. I am confident in this solution.", "hypothesis": "def f(a: int, b: int) {-}> int:\textbackslash{}n    \textbackslash{}"\textbackslash{}"\textbackslash{}"Returns 0 if both inputs are non{-}negative, the sum if both are negative, and the minimum if signs differ.\textbackslash{}"\textbackslash{}"\textbackslash{}"\textbackslash{}n    if a >= 0 and b >= 0:\textbackslash{}n        return 0\textbackslash{}n    elif a < 0 and b < 0:\textbackslash{}n        return a + b\textbackslash{}n    else:\textbackslash{}n        return a if a < b else b", "confidence\_score": 95\}} \\
            }
        }
    }
    & & \\ \\

    \theutterance \stepcounter{utterance}  
    & & & \multicolumn{2}{p{0.3\linewidth}}{
        \cellcolor[rgb]{0.9,0.9,0.9}{
            \makecell[{{p{\linewidth}}}]{
                \texttt{\tiny{[GM$|$GM]}}
                \texttt{Game Verdict: LOSS (turn limit reached)} \\
            }
        }
    }
    & & \\ \\

\end{supertabular}
}

\end{document}
